\chapter{Data model}
\label{ch:datamodel}
\begingroup
\sloppy

This chapter describes the classes that carry data between the tasks of ATTPCROOTv2: what each class
stores, which task fills it, under which branch name it is written, and how to read it back in a macro.
The persistent classes live in \code{AtData/} (detector and reconstruction data) and
\code{AtSimulationData/} (Monte Carlo truth); the PID container \code{AtPIDEvent} lives in
\code{AtReconstruction/}. Short reference tables with the same information are kept in
\url{docs/reference/data-model.md} and \url{docs/reference/branch-io-contracts.md}; where those files and
the source disagree, this chapter follows the source and says so.

\section{Overview: branches, containers and event objects}
\label{sec:dm-overview}

Every task registers its output with \code{FairRootManager::Register(name, folder, array, persistent)}.
The registered object is always a \code{TClonesArray}; the branch name and the class of the objects
inside the array are two different things. Two layouts occur:

\begin{itemize}
\item \textbf{Event branches.} The reconstruction branches (\code{AtRawEvent},
  \code{AtEventH}, \code{AtPatternEvent}, \code{AtTrackingEvent}, \code{AtPIDEvent}, \ldots) hold \emph{one} event
  object in slot 0. The tracks, hits and pads are members of that object, not separate
  \code{TClonesArray} entries. A task fills the slot with \code{ConstructedAt(0)} or placement
  \code{new (array[0])}.
\item \textbf{Point branches.} The simulation branches \code{AtTpcPoint} (\code{AtMCPoint}),
  \code{MCTrack} (\code{AtMCTrack}) and \code{AtSimulatedPoint} hold \emph{many} objects, one per
  Geant step, per particle or per electron cluster. The statement in
  \url{docs/reference/branch-io-contracts.md} that every branch holds one event object at slot 0 is
  true only for the event branches.
\end{itemize}

FairRoot writes the persistent branches of one run into a \code{TTree} named \code{cbmsim}; every macro
cited in this chapter opens that tree. The data flow of a full chain is

\begin{verbatim}
generators --(AtVertexPropagator)--> Geant (AtTpc)
     --> AtTpcPoint [AtMCPoint], MCTrack [AtMCTrack]
AtTpcPoint       --AtClusterizeTask--> AtSimulatedPoint
AtSimulatedPoint --AtPulseTask-------> AtRawEvent
  (experimental data: AtUnpackTask --> AtRawEvent)
AtRawEvent     --AtFilterTask--> AtRawEventFiltered (optional)
AtRawEvent     --AtPSAtask-----> AtEventH        [AtEvent: hits]
AtEventH       --AtPRAtask-----> AtPatternEvent  [AtTrack candidates]
AtPatternEvent --AtPIDTask-----> AtPIDEvent
AtPatternEvent --AtFitterTask--> AtTrackingEvent [AtFittedTrack]
                                 (+ AtFitMetadata)
\end{verbatim}

\section{Common base: AtBaseEvent}
\label{sec:dm-baseevent}

\code{AtBaseEvent} (\code{AtData/AtBaseEvent.h}, \code{ClassDefOverride(AtBaseEvent, 1)}) derives from
\code{TNamed} and is the base of \code{AtRawEvent}, \code{AtEvent}, \code{AtPatternEvent} (and through it
\code{AtFissionEvent}) and \code{AtTrackingEvent}. It holds the information that must follow an event
through every stage:

\begin{itemize}
\item \code{fEventID} (\code{GetEventID()}, \code{SetEventID()}); initialised to \code{-1}.
\item \code{fIsGood} (\code{IsGood()}, default \code{true}) and \code{fIsInGate}
  (\code{GetIsExtGate()}, default \code{false}).
\item \code{fTimestamp}, a \code{std::vector<ULong64\_t>} with one element by default; \code{GetTimestamp(i)}
  returns 0 for an index outside the vector, \code{SetNumberOfTimestamps(n)} resizes it.
\item \code{fAuxPadMap}, a \code{std::map<std::string, AtAuxPad>} of auxiliary channels (ion chamber,
  mesh, triggers) addressed by name: \code{AddAuxPad(name)}, \code{GetAuxPad(name)}, \code{GetAuxPads()}.
\item \code{fEventName}.
\end{itemize}

Every derived event has a constructor from \code{const AtBaseEvent\&}. \code{AtPSAtask} uses it to copy
the event ID, flags, timestamps and auxiliary pads from the raw event into the hit event
(\code{*event = *rawEvent} in \code{AtPSAtask::Exec}), which is how the ion-chamber trace and the event
number reach the later stages. \code{AtFitterTask::Exec} copies the event ID from the pattern event into
the tracking event explicitly; the comment there explains that without it a fit of a pre-gated file can
not be joined back to other per-event data (ion chamber) except by tree index.

\section{Pads and raw events}
\label{sec:dm-pads}

\subsection{AtPadBase and the pad augments}

\code{AtPadBase} (\code{TObject}, version 1) is an abstract base with a single pure virtual
\code{Clone()}. Two families derive from it: the pads themselves (\code{AtPad}, \code{AtAuxPad}) and
\emph{augments}, objects attached to a pad under a string key to carry extra per-pad information. The
augment classes present in \code{AtData/} are listed in Table~\ref{tab:dm-augments}.

\begin{table}[htbp]
\centering
\small
\begin{tabular}{p{3.4cm}p{5.4cm}p{5.6cm}}
\toprule
Class (version) & Content & Producer and key (from the source) \\
\midrule
\code{AtPadArray} (1) & \code{std::array<Double\_t,512>} trace & \code{AtPulse} key \code{"Q"} when
  charge saving is on; \code{AtPSADeconv} key \code{"Qreco"}; \code{AtPSAIterDeconv} key from
  \code{SetIterQName} (default \code{"Qreco"}) \\
\code{AtPadFFT} (1) & real and imaginary parts, $512/2+1$ points each & \code{AtFilterFFT},
  \code{AtPSADeconv}, \code{AtPSAIterDeconv}, key \code{"fft"}; \code{AtPSADeconv} also
  \code{"filter"} and \code{"Qreco-fft"} \\
\code{AtPadValue} (1) & one \code{Double\_t} & \code{AtGRAWUnpacker}, key \code{"lastCell"} \\
\code{AtPulserInfo} (1) & rise/fall begin, end and magnitude & \code{AtRemovePulser}, key
  \code{"pulserInfo"} \\
\code{AtElectronic\-Reference} (1) & an \code{AtPadReference} (CoBo, AsAd, AGET, channel) & no producer
  found by a search of the source \\
\code{AtGenericTrace} (1) & variable-length trace from another DAQ & stored in
  \code{AtRawEvent::fGTraceList}, not as an augment \\
\bottomrule
\end{tabular}
\caption{Classes derived from \code{AtPadBase} other than the pads. Keys are the strings passed to
\code{AtPad::AddAugment} in the files named.}
\label{tab:dm-augments}
\end{table}

\subsection{AtPad and AtAuxPad}

\code{AtPad} (version 3) holds one channel of the pad plane: the pad number \code{fPadNum} (the
reference number in \code{AtMap}), \code{fSizeID}, the pad centre \code{fPadCoord}, the flags
\code{fIsValid} and \code{fIsPedestalSubtracted}, the raw trace \code{fRawAdc}
(\code{std::array<Int\_t,512>}) and the calibrated trace \code{fAdc} (\code{std::array<Double\_t,512>}).
The trace length of 512 time buckets is fixed by the type. \code{GetADC()} returns \code{fAdc} even when
the pedestal has not been subtracted; it only writes a debug-level log message in that case.

Augments are kept in the member \code{fPadAugments}, a
\code{std::map<std::string, std::unique\_ptr<AtPadBase>>}, and are streamed with the pad. \code{AddAugment(name, ptr)} logs an error if the key is already used and then
replaces it; \code{ReplaceAugment} replaces silently. The templated
\code{GetAugment<T>(name)} performs the \code{dynamic\_cast} and returns \code{nullptr} if the key is
absent or the type differs:

\begin{verbatim}
const auto *q = pad->GetAugment<AtPadArray>("Q");   // nullptr if not saved
if (q) double v = q->GetArray(100);
\end{verbatim}

\code{AtAuxPad} derives from \code{AtPad}, adds \code{fAuxName}, and deletes \code{GetPadNum()},
\code{GetSizeID()} and \code{GetPadCoord()}. The header notes that these deletions are bypassed when the
object is used through an \code{AtPad} pointer or reference, because the functions are not virtual.

\subsection{AtRawEvent}

\code{AtRawEvent} (version 7) is the output of unpacking and of digitisation. Its members are
\begin{itemize}
\item \code{fPadList}, a \code{std::vector<std::unique\_ptr<AtPad>>}; add pads with \code{AddPad(...)}
  (perfect forwarding to an \code{AtPad} constructor, or a moved \code{unique\_ptr} of a derived type),
  read them with \code{GetPads()} or \code{GetPad(padNum)}. \code{GetPad} is a linear search over the
  list and returns \code{nullptr} for a pad that is absent;
\item \code{fFpnMap}, the fixed-pattern-noise channels keyed by \code{AtPadReference}
  (\code{AddFPN}, \code{GetFpn}, \code{GetFpnPads});
\item \code{fGTraceList}, a vector of \code{AtGenericTrace} for other acquisition systems;
\item \code{fSimMCPointMap}, a \code{std::multimap<Int\_t, std::size\_t>} from pad number to index in
  the \code{AtTpcPoint} array. This member is marked \code{//!} and is \emph{not} written to disk
  (Section~\ref{sec:dm-mctruth}).
\end{itemize}

\section{Hits and hit events}
\label{sec:dm-hits}

\subsection{AtHit}

\code{AtHit} (\code{TObject}, version 5) is a point in space with charge, produced by a PSA method from
one pad trace. Its members and accessors are:

\begin{itemize}
\item position \code{fPosition} (\code{ROOT::Math::XYZPoint}, \code{GetPosition()}) and its variance
  \code{fPositionVariance} (\code{GetPositionVariance()}, \code{GetPositionSigma()});
\item charge \code{fCharge} (\code{GetCharge()}) and \code{fChargeVariance};
  \code{fTraceIntegral} (default \code{-1}), the integral of the whole trace;
\item \code{fHitID} (\code{GetHitID()}), \code{fPadNum} (\code{GetPadNum()}) and \code{fHitMult}, the
  number of hits found in the same pad;
\item time: \code{fTimeStamp} (time bucket, \code{Int\_t}), \code{fTimeStampCorr} (centre of gravity)
  and \code{fTimeStampCorrInter} (interpolated);
\item Monte Carlo information: the member \code{fMCSimPointArray}, a vector of the nested
  struct \code{AtHit::MCSimPoint}, with \code{pointID}, \code{trackID}, \code{energy}, \code{eloss},
  \code{angle}, \code{A}, \code{Z}.
\end{itemize}

The hit itself has no cluster or track label; membership of a track is expressed by the
\code{AtTrack} that contains a copy of the hit.

\paragraph{Coordinate frame.} PSA writes the hit $z$ in the drift frame, measured from the pad plane.
The comment in \code{AtPSA::CalibrateZ} (\code{AtPSA.cxx}) states that the fitters convert with
$z_\mathrm{lab} = z_\mathrm{PadPlane} - z_\mathrm{digi}$, and that the optional Spyral-style calibration
(\code{fWindowTB > 0}) returns the opposite convention. Fitter outputs (vertex, smoothed positions) are
in the lab frame; compare them with hit positions only after this conversion.

\subsection{AtHitCluster and AtHitClusterFull}

\code{AtHitCluster} (version 4) derives from \code{AtHit} and represents a group of hits assumed to come
from the same charge deposit. It adds \code{fClusterID}, a $3\times3$ covariance matrix
\code{fCovMatrix} (\code{GetCovMatrix()}), the running weights used by \code{AddHit}, the charge-weighted
position \code{fPositionCharge} (\code{GetPositionCharge()}), \code{fLength} (default \code{-999}) and
\code{fClusterSize}. When a cluster is built with \code{AddHit(const AtHit\&)}, \code{fPosition} is
updated with weights charge $\times$ 1/variance per coordinate (\code{AtHitCluster::updatePosition}); the
header comment of \code{AtHitCluster.h} describes \code{fPosition} as unweighted, which does not match
the code.

The cluster builders in \code{AtTools/AtTrackTransformer.cxx}, which produce the clusters stored on
pattern-recognition tracks, do not call \code{AddHit}: they compute the position, charge and covariance
themselves and set them with \code{SetPosition}, \code{SetCharge} and \code{SetCovMatrix}. On those
clusters \code{GetClusterSize()} stays 0 and \code{GetPositionCharge()} stays at the origin.

\code{AtHitClusterFull} keeps a copy of every added hit and offers further covariance estimates. It has
a \code{ClassDefOverride(AtHitClusterFull, 1)} but no entry in \code{AtDataLinkDef.h}, so no dictionary
or streamer is generated for it and it can not be persisted.

\subsection{AtEvent}

\code{AtEvent} (version 6) is the hit-level event written by \code{AtPSAtask}. It holds
\code{fHitArray} (\code{std::vector<std::unique\_ptr<AtHit>>}), the total charge \code{fEventCharge}
(default \code{-100}), \code{fRhoVariance}, a pad multiplicity map, and the summed mesh signal \code{fMeshSig}, an
array of 512 \code{Float\_t} returned by \code{GetMesh()}.

\code{AddHit(...)} creates or moves a hit and, if its ID is still \code{-1}, sets the ID to the position
at which it was inserted. \code{AtPSA::Analyze} sorts the hits by time stamp after adding them
(\code{SortHitArrayTime()}), so in the output of PSA the hit ID is generally \emph{not} the position in
the vector. \code{GetHit(i)} takes a position; \code{GetHits()} returns the vector by const reference.
\code{GetHitArray()}, which returns a copy, is marked deprecated.

\section{Pattern recognition output}
\label{sec:dm-pattern}

\subsection{AtTrack}

\code{AtTrack} (\code{TObject}, version 5) is one candidate track. It contains
\begin{itemize}
\item \code{fTrackID} (default \code{-1}, see below);
\item \code{fHitArray}, owned copies of the hits (\code{std::vector<std::unique\_ptr<AtHit>>});
\item \code{fHitClusterArray}, a \code{std::vector<AtHitCluster>} of clusters built from those hits,
  which is what the fitters and the PID estimators consume;
\item \code{fPattern}, a \code{std::unique\_ptr<AtPatterns::AtPattern>} (line, ray, 2D circle, Y or
  fission pattern, all linked with \code{+} in \code{AtDataLinkDef.h});
\item geometric estimates from pattern recognition: \code{fGeoThetaAngle}, \code{fGeoPhiAngle},
  \code{fGeoRadius}, \code{fGeoCenter}, and \code{fChargeSign} ($+1/-1$ from the curvature, 0 unknown);
\item \code{fIsMerged}, \code{fVertexToZDist}, and the Bragg curve (\code{fBraggCurveValues} and the
  binned \code{BraggCurve} struct) filled by \code{AtPatternModification}.
\end{itemize}

Copying an \code{AtTrack} deep-copies the hit array. \code{SortClusterHitArrayZ()} sorts the clusters by
$z$; \code{AddClusterHit(std::shared\_ptr<AtHitCluster>)} moves the pointee into the vector.

\paragraph{Track ID semantics.} The meaning of \code{GetTrackID()} depends on the pattern-recognition
algorithm. \code{AtTrackFinderTC} and \code{AtTrackFinderHC} set it to the index of the cluster in the
clustering output (\code{track.SetTrackID(cluster\_index)}), skip empty clusters, and
\code{AtTrackFinderTC} can afterwards remove tracks in \code{SelectAndMergeTracks}. The ID is therefore a
cluster label and can differ from the position of the track in \code{GetTrackCand()}: an event with a
single track can have track ID 2. \code{AtTrackFinderHDBSCAN} uses a running counter,
\code{AtSampleConsensus} and \code{AtTrackFinderRiemann} the current number of tracks. Code that
connects fitted tracks or PID results to pattern tracks must match on the track ID, never on the array
index. The comments in \code{AtPIDTask.cxx} and in \code{AtTools/AtPID/AtSpyralPID.h} say the same.

\paragraph{No const access to the clusters.} The only accessor for the clusters is
\begin{verbatim}
// non-const member, returns a pointer
std::vector<AtHitCluster> *GetHitClusterArray();
\end{verbatim}
There is no const overload. Through a \code{const AtTrack\&}, for example in a loop
\code{for (const auto \&tr : patternEvent->GetTrackCand())}, the call does not compile; loop with a
non-const reference instead, as the macros in \code{macro/} do. The same holds for \code{GetMeanTime()},
\code{GetLinearRange()}, \code{GetFirstPoint()} and \code{GetLastPoint()}, which are also non-const.
\code{GetHitArray()} has both a const and a non-const overload.

\subsection{AtPatternEvent and AtFissionEvent}

\code{AtPatternEvent} (version 3) holds \code{std::vector<AtTrack> fTrackCand} (\code{GetTrackCand()},
const and non-const) and the hits not assigned to any track, \code{fNoise}. \code{GetNoiseHits()} is a
non-const member. \code{AtFissionEvent} (version 1) derives from \code{AtPatternEvent} and adds the beam
and fragment hits before space-charge correction and the correction strength \code{fLambda}; it is filled
by \code{AtFissionAnalysis} according to its header.

\section{Fit output}
\label{sec:dm-fit}

\subsection{AtTrackingEvent}

\code{AtTrackingEvent} (\code{ClassDef} version 1) is written by \code{AtFitterTask}. It stores
\begin{itemize}
\item \code{fFittedTrackArray}, a vector of \code{std::unique\_ptr<AtFittedTrack>}
  (\code{GetFittedTracks()});
\item \code{fTrackArray}, a copy of the \code{AtTrack} objects that were given to the fitter;
  \code{EventFit::AtFitter::FitEvent} sets it from the pattern event, and \code{AtGenfitter} with
  continuity merging sets it to the merged tracks. \code{GetTrackArray()} returns the vector
  \emph{by value}, so every call copies all tracks and their hits; keep the result in a local variable;
\item \code{fVertex}, \code{fVertexEnergy} (default \code{-10.0}) and \code{fGeoVertex}.
\end{itemize}
\code{AddFittedTrack} assigns the insertion index as track ID only if the fitted track still carries
\code{-1}. The fitters in the tree (\code{AtGenfitter}, \code{AtFitterUKF}) copy the ID of the source
\code{AtTrack}. A fitter can also refuse a track (\code{AtGenfitter} returns \code{nullptr} for a fitted
polar angle outside its window), so \code{GetFittedTracks()} is not parallel to the pattern tracks.

\subsection{AtFittedTrack}

\code{AtFittedTrack} (\code{ClassDef} version 3) is the result for one track. The per-particle members are
vectors indexed by a particle index (default 0), read with \code{.at()}, so asking for an index that
was never set throws \code{std::out\_of\_range}:

\begin{itemize}
\item \code{GetKinematics(i)} and \code{GetKinematicsXtr(i)}: struct \code{Kinematics} with
  \code{kineticEnergy}, \code{theta}, \code{phi} (all default \code{-1});
\item \code{GetParticleInfo(i)}: struct \code{ParticleInfo} with \code{idPDG} (\code{TString}),
  \code{charge}, \code{mass} (amu);
\item \code{GetVertex(i)}: \code{XYZVector}.
\end{itemize}

The two kinematics slots do \emph{not} mean the same thing for the two fitters in the tree:

\begin{table}[htbp]
\centering
\small
\begin{tabular}{p{2.4cm}p{6.0cm}p{6.0cm}}
\toprule
Fitter & \code{GetKinematics()} & \code{GetKinematicsXtr()} \\
\midrule
\code{AtGenfitter} & raw fit result at the first measurement point & back-extrapolated to the beam axis,
  including any manual energy-loss correction (the corrected value) \\
\code{AtFitterUKF} & at the vertex, after back-extrapolation (the corrected value) & at the first
  cluster, before back-extrapolation \\
\bottomrule
\end{tabular}
\caption{Meaning of the two kinematics slots, from \code{AtGenfitter.cxx} (the \code{ft->SetKinematics}
block) and \code{AtFitterUKF.cxx} (step 7, ``Build AtFittedTrack'').}
\label{tab:dm-kinslots}
\end{table}

The comment in \code{AtGenfitter.cxx} states that every consumer that wants the corrected energy must
read \code{GetKinematicsXtr()}. A macro that switches between the two fitters must switch slots as well.
For \code{AtGenfitter} the vertex and the smoothed positions are in the lab frame, in mm (converted from
genfit's cm in the code).

\code{GetTrackPropertiesStruct()} returns the \code{TrackProperties} struct: \code{initialPosition},
\code{initialPositionXtr}, \code{extrapolatedDistance}, \code{distancePOCA}, \code{trackLength},
\code{trackLengthXtr}, \code{estimateTotalCharge}, \code{estimateDeDx} (total charge divided by
\code{trackLengthXtr} when set through \code{SetTrackPropertiesStruct}), \code{trackPoints} and
\code{fSmoothedPositions} (\code{GetSmoothedPositions()}). \code{AtGenfitter} and \code{AtFitterUKF}
fill the smoothed positions; neither calls \code{SetTrackPropertiesStruct} in the code inspected, so the
remaining fields keep their defaults of \code{-1} for those fitters.

\code{GetTrackMetadata()} returns a reference to the \code{std::unique\_ptr<AtFitTrackMetadata>}, which
can be null. The class also keeps a set of members and methods marked \code{[[deprecated]]}
(\code{GetEnergyAngles}, \code{GetStats}, \code{GetTrackProperties}, \ldots) for the layout of earlier
versions; new code should use the accessors above.

\subsection{AtFitTrackMetadata and AtFitMetadata}

\code{AtFitTrackMetadata} (version 3) holds the statistics of one fit: \code{GetChi2()}, \code{GetNdf()},
\code{GetPValue()}, \code{GetFitConverged()}, \code{GetTrackID()}, and three flags:
\begin{itemize}
\item \code{GetGoodFit()}: the fit passed the fitter's quality cut. Tracks that fail are kept, not
  dropped; filter on this flag downstream. \code{AtGenfitter} sets it from convergence, $\chi^2/\mathrm{ndf}$
  and the kinetic-energy window, and always stores \code{PValue} as 0;
\item \code{GetMatEffects()}: material effects were active in the fit that produced the result;
\item \code{GetMatEffectsFallback()}: material effects were requested, that fit threw, and the result
  comes from the retry without material. The header comment recommends excluding these tracks before
  quoting a resolution.
\end{itemize}
\code{AtFitMetadata} (version 1) groups, per event, all metadata of all fits (a
\code{std::map<Int\_t, vector<unique\_ptr<AtFitTrackMetadata>>>} keyed by track ID). It is written only
if \code{AtFitterTask::SetFitMetadataBranch(name)} was called. \code{MCFitter::AtMCResult} (version 2)
derives from \code{AtFitTrackMetadata} and adds the sampled parameter map and the iteration number for
\code{AtMCFitterTask}. Unlike the other reconstruction branches, the \code{AtMCResult}, \code{SimEvent} and
\code{SimRawEvent} arrays hold one entry per stored Monte Carlo iteration
(\code{AtMCFitter::FillResultArrays}), not one event in slot 0.

\section{PID output: AtPIDEvent}
\label{sec:dm-pidevent}

\code{AtPIDEvent} (\code{TNamed}, \code{ClassDefOverride(AtPIDEvent, 3)}) is written by \code{AtPIDTask}
and holds four per-track vectors: \code{GetClassic()} (\code{AtTools::AtPIDResult}), \code{GetSpyral()}
(\code{AtTools::AtSpyralResult}), \code{GetHelix3D()} (\code{AtTools::AtHelix3DResult}) and
\code{GetCombined()} (\code{AtTools::AtPIDCombinedResult}). \code{GetDefault()} returns the Spyral
vector. The header documents the version history: version 2 added \code{fHelix3D} and version 3 added
\code{fCombined}; files written with an older version read back with the new vectors empty.
\code{fCombined} is filled only when \code{AtPIDTask::SetCombine(true)} is used, \code{fHelix3D} only
when the 3D estimator runs (default on). Each result struct carries the \code{trackID} of its source
track. In \code{AtSpyralResult} the member \code{failCode} is marked \code{//!} and is not on disk; the
comment explains that a streamed member would have changed the file layout. The PID chapter describes
the content of these structs.

\section{Simulation data}
\label{sec:dm-sim}

\subsection{AtMCPoint}

\code{AtMCPoint} (version 2) derives from \code{FairMCPoint}, which provides the track ID, position at
entrance (cm), momentum (GeV), time of flight (ns), length (cm) and energy loss (GeV), as documented in
the constructor comment. \code{AtMCPoint} adds the volume name (\code{GetVolName()}), copy number,
\code{GetEIni()} and \code{GetAIni()} (initial energy and angle of the particle, looked up in
\code{AtVertexPropagator} by \code{AtTpc::addHit}; zero for track ID 0), and \code{GetMassNum()},
\code{GetAtomicNum()}. Copy construction and assignment are deleted. The detector \code{AtTpc} registers
the collection as branch \code{AtTpcPoint}, persistent.

\code{AtStack::UpdateTrackIndex} rewrites the track ID of every point after the event so that it
indexes the stored \code{MCTrack} array. The comment in \code{AtTpc::reactionOccursHere} warns that track
ID 0 does not always mean the beam: in a reaction event no beam track is added and the first primary,
the heavy residual, has ID 0.

\subsection{AtMCTrack and AtStack}

\code{AtStack} (registered branch \code{MCTrack}, persistent) stores one \code{AtMCTrack} per kept
particle with PDG code, mother ID, start momentum, start position and time, and number of points.
Primaries are always kept; secondaries are filtered by \code{StoreSecondaries} (default \code{kTRUE}),
\code{SetMinPoints} (default 1), \code{SetEnergyCut} (default 0) and \code{StoreMothers} (default
\code{kTRUE}, keeps the ancestors of every kept track), as set in the \code{AtStack} constructor.

\subsection{AtSimulatedPoint}

\code{AtSimulatedPoint} (\code{AtDigitization/}, version 3) is the intermediate object between Geant
points and pad traces: one electron cluster with the index of its \code{AtMCPoint}, an event ID, a
cluster ID, a charge and a position in (mm, mm, \textmu s) according to the member comment.
\code{AtClusterizeTask} creates the array with the class name returned by the clusterizer
(\code{AtSimulatedPoint}, or \code{AtSimulatedLine} for \code{AtClusterizeLine}).

\subsection{AtVertexPropagator}

\code{AtVertexPropagator} is a singleton (\code{AtVertexPropagator::Instance()}) that carries state
between the generators of one event and between generation and transport. It is never written to a
branch, although it is linked with \code{+} and has \code{ClassDef} version 2. It holds
\begin{itemize}
\item the reaction vertex (\code{GetVx/Vy/Vz}), the beam entrance point (\code{GetInVx/...}), the beam
  momentum and energy at the vertex (\code{GetPx/Py/Pz}, \code{GetEnergy}), set by
  \code{AtTpc::startReactionEvent} when the beam has lost the randomly drawn energy
  \code{GetRndELoss()};
\item the beam mass, nominal beam energy, $A$ and $Z$;
\item per-track energy and angle maps (\code{SetTrackEnergy(trackID, E)}, \code{GetTrackEnergy}),
  filled by the reaction generators (\code{AtTPC2Body}, \code{AtTPCXSReader}, \code{AtTPC\_d2He},
  \ldots) and read by \code{AtTpc::addHit} and by decay generators;
\item the momenta and excitation energies of scattered and recoil particles, validity flags and the
  $d$--$^2$He vertex.
\end{itemize}
The flag \code{IsBeamEvent()} alternates: \code{AtReactionGenerator::ReadEvent} calls \code{EndEvent()},
which inverts it, when the generator is flagged as the final one. The comment in \code{AtTpc.cxx} warns
that, because the flag is flipped during generation, its sense is inverted by the time Geant steps.

\section{Monte Carlo truth on hits}
\label{sec:dm-mctruth}

The \code{AtHit::MCSimPoint} entries are filled only when all of the following hold:
\begin{enumerate}
\item \code{AtPulseTask::SetSaveMCInfo()} was called (default off), so that the task fills the pad to
  point map and stores it with \code{AtRawEvent::SetSimMCPointMap};
\item \code{AtPSAtask} runs in the \emph{same} \code{FairRunAna} as \code{AtPulseTask}, because
  \code{fSimMCPointMap} is transient (\code{//!}) and is lost when the raw event is written;
\item the \code{AtTpcPoint} array is available to \code{AtPSAtask} under the branch name set by
  \code{SetSimlulatedPointBranch} (default \code{"AtTpcPoint"}; note the spelling of the setter).
\end{enumerate}
\code{AtPSA::TrackMCPoints} then adds one \code{MCSimPoint} per \code{AtMCPoint} that deposited charge
in the pad of the hit, without checking the time bucket.

\section{Branch input/output of the standard tasks}
\label{sec:dm-branches}

Table~\ref{tab:dm-branches} lists the default branch names and persistence of the standard tasks,
taken from the constructors, member initialisers and \code{Init()} methods. All defaults can be changed
with the setters named in the last column, except where the column says ``fixed''.

\begin{table}[htbp]
\centering
\footnotesize
\begin{tabular}{p{3.3cm}p{3.0cm}p{4.0cm}p{4.0cm}}
\toprule
Task & Input branch & Output branch [class] & Persist default; setters \\
\midrule
\code{AtTpc} (detector) & -- & \code{AtTpcPoint} [\code{AtMCPoint}, many] & true; fixed \\
\code{AtStack} & -- & \code{MCTrack} [\code{AtMCTrack}, many] & true; fixed \\
\code{AtClusterizeTask} & \code{AtTpcPoint} & \code{AtSimulatedPoint} [\code{AtSimulatedPoint},
  many] & false; names fixed, \code{SetPersistence} \\
\code{AtPulseTask} & \code{AtSimulated\-Point}, \code{AtTpcPoint} & \code{AtRawEvent}
  [\code{AtRawEvent}]; re-registers \code{AtTpcPoint} & true (\code{AtTpcPoint}: false);
  \code{SetOutputBranch}, \code{SetPersistence}, \code{SetPersistence\-AtTpcPoint} \\
\code{AtUnpackTask} & external file & \code{AtRawEvent} [\code{AtRawEvent}] & true;
  \code{SetOuputBranchName}, \code{SetPersistence} \\
\code{AtFilterTask} & \code{AtRawEvent} & \code{AtRawEventFiltered} [\code{AtRawEvent}] & false;
  \code{SetInputBranch}, \code{SetOutputBranch} \\
\code{AtPSAtask} & \code{AtRawEvent} (+ \code{AtTpcPoint} if present) & \code{AtEventH}
  [\code{AtEvent}] & false; \code{SetInputBranch}, \code{SetOutputBranch} \\
\code{AtDataCleaning\-Task} & \code{AtEventH} & \code{AtEventCleaned} [\code{AtEvent}] & true;
  \code{SetInputBranch}, \code{SetOutputBranch} \\
\code{AtSpaceCharge\-CorrectionTask} & \code{AtEventH} & \code{AtEventCorrected} [\code{AtEvent}] &
  true; \code{SetInputBranch}, \code{SetOutputBranch} \\
\code{AtPRAtask} & \code{AtEventH} & \code{AtPatternEvent} [\code{AtPatternEvent}] & false;
  \code{SetInputBranch}, \code{SetOutputBranch} \\
\code{AtSampleConsensus\-Task} & \code{AtEventH} & \code{AtPatternEvent} [\code{AtPatternEvent}] &
  false; \code{SetInputBranch}, \code{SetOutputBranch} \\
\code{AtRansacTask} & \code{AtEventH} & \code{AtPatternEvent} [\code{AtPatternEvent}] & false;
  \code{SetInputBranchName}, \code{SetOutputBranchName} \\
\code{AtPattern\-ModificationTask} & \code{AtPatternEvent} (+ \code{AtRawEvent}, \code{AtEvent}) &
  \code{AtPatternEventModified} [\code{AtPatternEvent}] & false; \code{SetPersistence} \\
\code{AtPIDTask} & \code{AtPatternEvent} & \code{AtPIDEvent} [\code{AtPIDEvent}] & true;
  \code{SetInputBranch}, \code{SetOutputBranch} \\
\code{AtFitterTask} & \code{AtPatternEvent} (+ optional raw/hit event) & \code{AtTrackingEvent}
  [\code{AtTrackingEvent}]; optional metadata [\code{AtFitMetadata}] & false; \code{SetInputBranch},
  \code{SetOutputBranch}, \code{SetRawEventBranch}, \code{SetEventBranch},
  \code{SetFitMetadataBranch} \\
\code{AtMCFitterTask} & \code{AtPatternEvent} & \code{AtMCResult} [\code{AtMCResult}, many],
  \code{SimEvent} [\code{AtEvent}, many], \code{SimRawEvent} [\code{AtRawEvent}, many] & true, false, false;
  \code{SetPatternBranchName}, \code{SetSave*} \\
\bottomrule
\end{tabular}
\caption{Default branches of the standard tasks. ``many'' marks point branches with one object per
entry of the \code{TClonesArray}; all other branches hold one event in slot 0.}
\label{tab:dm-branches}
\end{table}

Three consequences of the defaults are worth stating:
\begin{itemize}
\item \code{AtPSAtask}, \code{AtPRAtask} and \code{AtFitterTask} are \emph{not} persistent by default.
  A reconstruction macro must call \code{SetPersistence(kTRUE)} on every task whose output it wants on
  disk.
\item \code{AtPRAtask}, \code{AtSampleConsensusTask} and \code{AtRansacTask} write the same default
  branch; only one can run per chain unless the names are changed.
\item The fitter task reads \code{AtPatternEvent} by default. To fit the output of
  \code{AtPatternModificationTask}, set \code{SetInputBranch("AtPatternEventModified")}.
\end{itemize}

\section{Empty events and repeated entries}
\label{sec:dm-repeat}

A task that returns from \code{Exec} before clearing its output array leaves the object of the previous
event in the array, and FairRoot writes that object again. Whether this happens depends on the task:

\begin{itemize}
\item \code{AtPRAtask::Exec} calls \code{fPatternEventArray.Delete()} first and creates an
  \code{AtPatternEvent} only if the number of hits $n$ satisfies $10 < n < 5000$ (\code{fMinNumHits},
  \code{fMaxNumHits}) and the algorithm returned an event. Otherwise the pattern array is empty for that
  event.
\item \code{AtPIDTask::Exec} also calls \code{Delete()} first, so for such an event the \code{AtPIDEvent}
  array is empty.
\item \code{AtFitterTask::Exec} returns immediately when the pattern array is empty, \emph{before}
  \code{fTrackingEventArray.Delete()}. The \code{AtTrackingEvent} of the previous event therefore stays
  in slot 0 and is written again in the current entry, including its event ID.
\item \code{AtPSAtask::Exec} returns before touching its output when the raw-event array is empty, with
  the same effect on \code{AtEventH}.
\end{itemize}

In fit output trees this shows up as an \code{AtTrackingEvent} repeated in consecutive entries, the
repeated copy sitting next to an empty \code{AtPIDEvent} array. A macro that loops over
\code{GetFittedTracks()} of every entry counts those tracks twice. Two protections work: skip entries
whose \code{AtPIDEvent} array is empty, or compare the event ID of the tracking event with the one of the
previous entry. The macro in Section~\ref{sec:dm-reading} uses the first.

\section{Reading the output in a macro}
\label{sec:dm-reading}

Open the file, take the \code{cbmsim} tree, give each branch a \code{TClonesArray} pointer, and cast slot
0 to the event class. The libraries must be loaded first (\code{gSystem->Load("libAtReconstruction.so")}
loads its dependencies, among them the \code{AtData} dictionary). A minimal hit loop, from
\url{macro/Unpack_HDF5/a1975/D2_UKF/_archive/ckhits.C} (reformatted):

\begin{verbatim}
auto f = TFile::Open("/mnt/f/a1975/reco_d2/run_0300_psa_" + tag + ".root");
auto t = (TTree *)f->Get("cbmsim");
TClonesArray *ar = nullptr;
t->SetBranchAddress("AtEventH", &ar);
for (long i = 0; i < n; i++) {
   t->GetEntry(i);
   auto e = (AtEvent *)ar->At(0);
   if (e) hpe.push_back(e->GetHits().size());
}
\end{verbatim}

A track and cluster loop over \code{AtPatternEvent}, from
\url{macro/Unpack_HDF5/a1975/UKF/pipeline/praDump_a1975.C}. Note the non-const loop variable required by
\code{GetHitClusterArray()}:

\begin{verbatim}
auto *arr = new TClonesArray("AtPatternEvent");
t->SetBranchAddress("AtPatternEvent", &arr);
for (Long64_t i = 0; i < t->GetEntries(); ++i) {
   t->GetEntry(i);
   auto *pe = (AtPatternEvent *)arr->At(0);
   if (pe)
      for (auto &tr : pe->GetTrackCand()) {
         auto *cl = tr.GetHitClusterArray();
         nClus += cl->size();
         for (auto &c : *cl) {
            auto p = c.GetPosition();
            cks += p.X() + p.Y() + p.Z();
         }
         rsum += tr.GetGeoRadius();
      }
}
\end{verbatim}

Fitted tracks joined to PID results by track ID, from
\url{macro/Unpack_HDF5/a1975/D2_UKF/ex_dt_a1975.C} (lines 126--201, shortened). The test on the
entry counts of both arrays skips the repeated entries of Section~\ref{sec:dm-repeat}, and the map
implements the match on track ID:

\begin{verbatim}
TTree *tu = (TTree *)fu->Get("cbmsim");
TClonesArray *te = nullptr, *pe = nullptr;
tu->SetBranchAddress("AtTrackingEvent", &te);
tu->SetBranchAddress("AtPIDEvent", &pe);
for (Long64_t i = 0; i < N; ++i) {
   tu->GetEntry(i);
   if (pe->GetEntries() == 0 || te->GetEntries() == 0)
      continue;
   auto *pidev = (AtPIDEvent *)pe->At(0);
   auto *ev = (AtTrackingEvent *)te->At(0);
   std::map<int, AtFittedTrack *> fmap;
   for (auto &ft : ev->GetFittedTracks())
      if (ft) fmap[ft->GetTrackID()] = ft.get();
   for (auto &sr : pidev->GetSpyral()) {
      if (!sr.valid) continue;
      auto it = fmap.find(sr.trackID);
      if (it == fmap.end()) continue;
      auto *ft = it->second;
      auto &k = ft->GetKinematicsXtr(); // genfit: back-extrapolated
      double ndf = ft->GetTrackMetadata()->GetNdf();
      double chi2 = ft->GetTrackMetadata()->GetChi2();
      ...
\end{verbatim}

When reading the fitter output, decide which kinematics slot to use from
Table~\ref{tab:dm-kinslots}, and check \code{GetTrackMetadata()} for \code{nullptr} if the file may come
from a fitter that does not set it.

\section{LinkDef conventions and persistence}
\label{sec:dm-linkdef}

Each library has a \code{*LinkDef.h} from which rootcling generates the dictionary. The suffix of each
entry decides what is generated:
\begin{itemize}
\item \code{\#pragma link C++ class X + ;} generates the dictionary with a full I/O streamer. This is
  needed for any class written to a file, directly or as a member of a written class.
\item \code{\#pragma link C++ class X - ! ;} generates reflection only: \code{-} suppresses the streamer
  and \code{!} the \code{operator>>}.
\end{itemize}
The project rule in \url{.claude/CLAUDE.md} and \url{docs/contributing/code-style.md} is to use
\code{+} only for classes written to disk and \code{-!} for tasks, algorithms and models, and to use the
ROOT typedefs (\code{Int\_t}, \code{Double\_t}) only in persisted classes. The existing code does not
follow the rule everywhere: many tasks are linked with \code{+}, for example \code{AtPSAtask},
\code{AtPRAtask}, \code{AtFitterTask} and \code{AtPIDTask} in \code{AtReconstructionLinkDef.h} and
\code{AtPulseTask}, \code{AtClusterizeTask} and \code{AtUnpackTask} in their libraries. \code{CLAUDE.md}
cites \code{AtFitterTask} as an example of \code{-!}, which does not match the LinkDef. Apply the rule
to new classes.

Inside a persisted class, a member marked with \code{//!} in its comment is transient: it is not
streamed and has its default value after reading. The transient members relevant to analysis are
\code{AtRawEvent::fSimMCPointMap} and \code{AtSpyralResult::failCode}. Every class that is to be
persisted must also appear in the LinkDef; \code{AtHitClusterFull} has a \code{ClassDef} but no LinkDef
entry and therefore cannot be written.

The result structs stored in \code{AtPIDEvent} are linked as structs with \code{+} in
\code{AtTools/AtToolsLinkDef.h}, and the vectors of three of them in
\code{AtReconstructionLinkDef.h}. \code{AtData} links \code{std::vector<ROOT::Math::XYZPoint>} for the
smoothed positions.

\subsection{Schema evolution}

No \code{\#pragma read} rules exist in the LinkDef files; the classes rely on ROOT's automatic schema
evolution, in which a member missing from an older file keeps the value set by the default
constructor. The project rule (\code{docs/contributing/code-style.md}) is to increment the
\code{ClassDef} version whenever the member layout of a streamed class changes. The current versions of
the main classes are given in Table~\ref{tab:dm-versions}, with the changes that the code comments or
the git history document.

\begin{table}[htbp]
\centering
\small
\begin{tabular}{p{3.4cm}p{1.2cm}p{9.8cm}}
\toprule
Class & Version & Documented changes \\
\midrule
\code{AtRawEvent} & 7 & -- \\
\code{AtEvent} & 6 & -- \\
\code{AtHit} & 5 & -- \\
\code{AtHitCluster} & 4 & -- \\
\code{AtPad} & 3 & -- \\
\code{AtPatternEvent} & 3 & -- \\
\code{AtTrack} & 5 & 4 $\to$ 5 added \code{fChargeSign} (commit 1fa0f5ae) \\
\code{AtTrackingEvent} & 1 & -- \\
\code{AtFittedTrack} & 3 & 2 $\to$ 3 added \code{fSmoothedPositions} to \code{TrackProperties} (commit
  8927b00a) \\
\code{AtFitTrackMetadata} & 3 & 1 $\to$ 2 added \code{fGoodFit} (273b9b1e); 2 $\to$ 3 added the
  material-effects flags (a2b14d8b) \\
\code{AtPIDEvent} & 3 & 2 added \code{fHelix3D}; 3 added \code{fCombined} (header comment) \\
\code{AtMCPoint} & 2 & -- \\
\code{AtMCResult} & 2 & -- \\
\bottomrule
\end{tabular}
\caption{\code{ClassDef} versions in the current source. A dash means that no change is documented in
the code comments; earlier history was not traced.}
\label{tab:dm-versions}
\end{table}

The classes \code{AtTrackingEventOld}, \code{AtFittedTrackOld} (declared \code{[[deprecated]]}) and
\code{MCFitter::AtMCResultOld} are still linked with \code{+}. They are the layouts used before the
fitter refactoring (commit 3d7023d3, according to the git log of \code{AtFittedTrackOld.h}), and are
used by \code{AtFitterTaskOld} and \code{AtMCFitterTaskOld}; keep them to read files written with that
code.

\section{Helpers: AtContainerManip and the data observer}
\label{sec:dm-helpers}

\code{AtData/AtContainerManip.h} (namespace \code{ContainerManip}) contains templates used throughout the
code to move between containers of objects and of \code{unique\_ptr}: \code{GetObjectVector} (copies,
used by \code{AtTrack::GetHitArrayObject()}), \code{GetPointerVector}, \code{GetConstPointerVector},
\code{MoveFromVector}, \code{GetMedian}, and \code{SetHistFromData}/\code{CreateHistFromData} to put an
array into a histogram. In the current source \code{CreateHistFromData} ignores its \code{xMin} and
\code{xMax} arguments and always uses the range $[0, n]$.

\code{AtDataSubject.h} and \code{AtDataObserver.h} define the observer pattern in namespace
\code{DataHandling}: an \code{AtSubject} keeps a set of \code{AtObserver} pointers and calls
\code{Notify()} on change; \code{AtSimpleType<T>} wraps a value and notifies only when it changes. These
classes are not persisted; their users are the event-display classes in \code{AtEventDisplay/}
(\code{AtViewerManager}, the tabs).

\section{Pitfalls}
\label{sec:dm-pitfalls}

\begin{itemize}
\item \textbf{Track ID is not an index.} With \code{AtTrackFinderTC} and \code{AtTrackFinderHC} the
  track ID is a cluster label. Join fitted tracks, PID results and pattern tracks on
  \code{GetTrackID()} / \code{trackID}.
\item \textbf{Repeated tracking events.} \code{AtFitterTask} rewrites the previous
  \code{AtTrackingEvent} for an event without pattern output. Skip entries with an empty
  \code{AtPIDEvent} array, or check the event ID.
\item \textbf{Kinematics slots.} \code{GetKinematics()} is raw for \code{AtGenfitter} and corrected for
  \code{AtFitterUKF}; \code{GetKinematicsXtr()} is the opposite.
\item \textbf{No const cluster accessor.} \code{AtTrack::GetHitClusterArray()} is non-const; loop over
  tracks with a non-const reference.
\item \textbf{By-value getters.} \code{AtTrackingEvent::GetTrackArray()} and \code{AtEvent::GetHitArray()}
  return copies.
\item \textbf{Hit ID after PSA.} Hits are sorted by time after insertion; the hit ID is not the position
  in \code{GetHits()}.
\item \textbf{Cluster bookkeeping.} Clusters made by \code{AtTrackTransformer} report
  \code{GetClusterSize() == 0}.
\item \textbf{Persistence defaults.} PSA, pattern recognition and fit tasks do not write their output
  unless \code{SetPersistence(kTRUE)} is called.
\item \textbf{MC truth on hits} needs \code{SetSaveMCInfo()} and PSA in the same run as digitisation.
\item \textbf{Frames.} Hit $z$ is in the drift frame; fitter vertices and smoothed positions are in the
  lab frame.
\end{itemize}
\endgroup
